\section{二进制炸弹}

\subsection{实验概述}

介绍本次实验的目的意义、目标、要求及安排等。正式撰写时，应依据 \path{experiments/lab_binary_bomb/requirement_summary.md} 补充，不得脱离原始实验要求。

\subsection{实验内容}

介绍本次实验的主要内容。如本实验分为多个阶段，应使用如下结构展开：

\subsubsection{阶段1：待填写阶段名称}

\begin{enumerate}[label=\arabic*.]
    \item \textbf{任务描述：}说明本阶段需要完成的具体任务。
    \item \textbf{实验设计：}说明解题思路、技术路线和拟采用的方法。
    \item \textbf{实验过程：}依据 \path{experiments/lab_binary_bomb/evidence/commands.md}、\path{experiments/lab_binary_bomb/outputs/} 和截图索引描述真实操作过程。
    \item \textbf{实验结果：}给出本阶段实验结果和必要分析，关键结论必须能追溯到证据编号。
\end{enumerate}

\subsection{实验设计}

给出解题思路分析和拟采用的技术、方法、工具链和关键数据结构。若涉及 GDB、反汇编、栈帧、ELF 或 C/ASM 接口，应结合图表和证据说明。

\subsection{实验过程}

分步骤说明实验的具体操作过程。每个关键步骤应包含：目的、命令或操作、输出位置、截图编号或证据编号。

\subsection{实验结果}

给出实验结果和必要的结果分析。不得只贴截图，应解释截图或输出说明了什么。

\subsection{实验小结}

对本次实验使用的理论、技术、方法和结果进行总结。可从 \path{journals/experiment_feelings_and_reflections.md} 中提取个人体会，但不得编造未发生的实验经历。
